\subsection{Irreducible Quadratic Factors} 
Suppose that we've discovered an imaginary root of a polynomial. How can we use it?
\begin{Example}
The polynomial $x^4-4x^3+3x^2+2x-6$ has $1-i$ as one of its roots. Use this fact to factor it fully.
\begin{lpic}{}{5}
\end{lpic}
... That's... not good. There's got to be an easier way.
\end{Example}
If we know one complex root, its \makeblank is also a root. We can put these together to get an \emph{irreducible quadratic factor}.
\begin{lpic}{}{5}
\twocol{-1}{5}{Method 1: Think of Factors}{Method 2: Think of Roots}
\end{lpic}
Okay. Let's try that again.
\begin{Example}
The polynomial $x^4-4x^3+3x^2+2x-6$ has $1-i$ as one of its roots. Use this fact to factor it fully.
\begin{lpic}{}{11}
\regtextleft{1}{-1}{We have irreducible polynomial factor:};
\regtextleft{1}{-11}{$x^4-4x^3+3x^2+2x-6={}$};
\end{lpic}
\end{Example}